\documentclass[12pt,a4paper]{ctexart}
\usepackage[margin=2.2cm]{geometry}
\usepackage{graphicx}
\usepackage{booktabs}
\usepackage{amsmath}
\usepackage{float}
\usepackage{hyperref}
\usepackage{xcolor}
\usepackage{listings}

\lstset{
  language=Python,
  basicstyle=\ttfamily\small,
  keywordstyle=\color{blue!70!black},
  commentstyle=\color{green!40!black},
  stringstyle=\color{orange!60!black},
  showstringspaces=false,
  breaklines=true,
  frame=single,
  tabsize=4
}

\newcommand{\reporttitle}{大实验报告 —— AgeNet (APPA-REAL) 年龄预测}
\newcommand{\subtask}{终端训练与评估}
\newcommand{\authname}{姓名}
\newcommand{\authid}{学号}

\title{\reporttitle}
\author{\authname：\underline{李祥熙} \quad \authid：\underline{2234412799}\\
        \authname：\underline{储正非} \quad \authid：\underline{2233412534}\\
        \authname：\underline{袁宇霄} \quad \authid：\underline{2235015935}\\
        }
\date{\today}

\begin{document}
\maketitle
\tableofcontents
\newpage

\section{实验目的}

本实验完成一个基于 APPA-REAL 人脸数据集的端到端年龄预测系统。核心任务是在 PyTorch 中实现指定的 \texttt{AgeNet(nn.Module)}，输入 \texttt{128$\times$128} RGB 人脸图像，输出 0 到 100 岁共 101 个年龄类别。

本次实验的主要目标如下：
\begin{enumerate}
    \item 理解年龄估计任务中分类建模与回归评价相结合的思想；
    \item 使用 PyTorch 实现课程要求的 AgeNet 卷积神经网络；
    \item 完成 APPA-REAL 图像数据读取、预处理、数据增强、训练、验证和测试预测；
    \item 使用平均绝对误差 MAE 评价年龄估计结果；
    \item 保存验证集 MAE 最优的模型权重，并输出测试集预测年龄文件；
    \item 根据训练日志分析模型收敛过程和实验结果。
\end{enumerate}

\section{实验过程}
\subsection{实验环境}
\begin{itemize}
  \item 操作系统：Linux
  \item 编程语言：Python 3
  \item 深度学习框架：PyTorch
  \item 关键依赖：torchvision, pandas, Pillow, OpenCV, imgaug, numpy
  \item 计算资源：优先使用 GPU（训练日志显示使用了 CUDA）
\end{itemize}


\subsection{年龄估计建模方法}

年龄估计既可以直接建模为回归问题，也可以转化为分类问题。本实验采用分类训练、期望回归的方式：模型输出 101 维 logits，每一维对应一个年龄类别。训练阶段将年龄四舍五入并裁剪到 0--100 范围，使用交叉熵损失进行监督；预测阶段对 logits 做 softmax 得到年龄概率分布，再计算概率分布的期望作为最终预测年龄。

设网络输出为 $f(x)$，年龄类别为 $a \in \{0,1,\ldots,100\}$，则预测过程为：
\[
P(a|x)=\mathrm{softmax}(f(x)),
\qquad
\hat{y}=\sum_{a=0}^{100} a \cdot P(a|x).
\]

这种方法避免了只取最大概率类别时的离散跳变，也比直接回归单个年龄值更容易利用分类损失稳定训练。由于 APPA-REAL 的年龄标签来自多名标注者的表观年龄估计，标签本身具有一定主观性，使用概率分布表达预测结果更符合任务特点。

\subsection{评价指标}

实验主要使用平均绝对误差 MAE 评价模型：
\[
\mathrm{MAE}=\frac{1}{N}\sum_{i=1}^{N}|\hat{y}_i-y_i|.
\]

MAE 的单位是岁，含义是模型平均预测偏差，因此比分类准确率更适合年龄估计任务。训练过程中每轮在验证集上计算 MAE，并以验证集 MAE 最小时的参数作为最终模型。

\subsection{数据集与预处理}
\subsubsection{数据集}
本实验使用 APPA-REAL 数据集（人脸裁剪后的 `\_face.jpg` 图像与相应 CSV 标注）。关键点：
\begin{itemize}
  \item 输入尺寸：128 $\times$ 128 RGB；在数据加载器中通过 \texttt{cv2.resize(..., (128,128))} 强制为 128 分辨率；
  \item 标签：每张图像的"apparent age"（取整并裁剪到 0--100）；训练时通过对真实标签加上带标准差的噪声进行数据增强（仅训练集启用）；
  \item 归一化：在 \texttt{FaceDataset} 中使用 transforms.Normalize(mean=[0.5,0.5,0.5], std=[0.5,0.5,0.5])；
  \item 批量大小与并行：默认 \texttt{BATCH\_SIZE=12}, \texttt{NUM\_WORKER=8}。
\end{itemize}

项目要求将数据集整理为如下目录结构：

\begin{lstlisting}[caption={APPA-REAL 数据目录结构}]
DATASET/
|-- appa-real-release/
|   |-- gt_avg_train.csv
|   |-- gt_avg_valid.csv
|   |-- gt_avg_test.csv
|   |-- train/
|   |-- valid/
|   |-- test/
\end{lstlisting}

其中训练集和验证集 CSV 文件提供可用于监督学习的年龄标签；测试集没有公开真实年龄标签，因此训练结束后只能生成测试集预测文件，不能直接计算测试 MAE。

项目中的 \texttt{helperT.py} 负责读取数据。程序从训练、验证和测试三个 CSV 文件中读取图片文件名，再到对应的 \texttt{train}、\texttt{valid} 和 \texttt{test} 目录中查找 \texttt{\_face.jpg} 人脸裁剪图。

项目还使用 \texttt{ignore\_list.csv} 跳过部分异常样本，避免缺失图片或异常标注影响训练流程。

\subsubsection{基础预处理}

输入图片统一缩放为 \texttt{128$\times$128}，并完成如下处理：
\begin{enumerate}
    \item 使用 OpenCV 读取人脸图片；
    \item 将图片 resize 到 \texttt{128$\times$128}；
    \item 将 BGR 格式转换为 RGB 格式；
    \item 使用 \texttt{ToTensor} 转换为 PyTorch 张量；
    \item 使用均值 \texttt{[0.5, 0.5, 0.5]} 和标准差 \texttt{[0.5, 0.5, 0.5]} 归一化，使输入大致落在 \texttt{[-1, 1]}。
\end{enumerate}

\subsubsection{训练集增强}

为了提高模型对光照、姿态和成像质量变化的鲁棒性，训练集使用了多种数据增强：
\begin{enumerate}
    \item 加性高斯噪声；
    \item 高斯模糊；
    \item 随机旋转，范围约为 $-20^\circ$ 到 $20^\circ$；
    \item 小幅缩放和平移；
    \item Gamma 对比度变化；
    \item 随机水平翻转。
\end{enumerate}

训练时还会根据 \texttt{apparent\_age\_std} 对年龄标签加入随机扰动，然后四舍五入并裁剪到 0--100 范围。这一处理利用了 APPA-REAL 标注标准差所反映的不确定性，有助于降低模型对单一整数标签的过拟合。

\subsection{实现步骤}

\subsubsection{实现 AgeNet}

AgeNet 是一个 VGG 风格的卷积神经网络，包含 5 个卷积特征提取阶段和 3 个全连接层。网络输入为 $3 \times 128 \times 128$，每经过一次 $2\times2$ 最大池化，空间尺寸减半，最终得到 $512 \times 4 \times 4$ 的特征图。

\begin{center}
\texttt{128 -> 64 -> 32 -> 16 -> 8 -> 4}
\end{center}

展平后的 8192 维特征依次经过两个 4096 维全连接层和一个 101 维输出层，输出 0 到 100 岁的类别 logits。

\begin{table}[!htbp]
    \centering
    \caption{AgeNet 网络结构}
    \begin{tabular}{p{2.2cm}p{7.0cm}p{3.8cm}}
        \toprule
        模块 & 层结构 & 输出尺寸变化 \\
        \midrule
        Block 1 & \texttt{Conv 7x7, 3->64} + ReLU + MaxPool & $128\times128 \to 64\times64$ \\
        Block 2 & \texttt{Conv 5x5, 64->128} + ReLU + \texttt{Conv 5x5, 128->128} + ReLU + MaxPool & $64\times64 \to 32\times32$ \\
        Block 3 & 3 个 \texttt{3x3} 卷积，通道 \texttt{128->256->256->256} + MaxPool & $32\times32 \to 16\times16$ \\
        Block 4 & 3 个 \texttt{3x3} 卷积，通道 \texttt{256->512->512->512} + MaxPool & $16\times16 \to 8\times8$ \\
        Block 5 & 3 个 \texttt{3x3} 卷积，通道 \texttt{512->512->512->512} + MaxPool & $8\times8 \to 4\times4$ \\
        分类器 & Flatten + \texttt{8192->4096->4096->101} & 输出 101 类 \\
        \bottomrule
    \end{tabular}
\end{table}

\subsubsection{数据集加载、训练和验证}

\begin{enumerate}
  \item 编写数据加载器 \texttt{FaceDataset} 与增强管线（使用 \texttt{imgaug}）；
  \item 在 \texttt{train\_ageNet.py} 中构建训练循环（100 epochs）、使用 SGD 优化器、交叉熵损失，并在每个 epoch 后计算验证 MAE（以概率分布的期望作为年龄预测）；训练参数如表 \ref{parameter}.

    \begin{table}[h!]
    \centering
    \label{parameter}
    \caption{训练参数设置}
    \begin{tabular}{ll}
        \toprule
        参数 & 设置 \\
        \midrule
        训练轮数 & 100 epochs \\
        Batch size & 12 \\
        优化器 & SGD \\
        学习率 & 0.001 \\
        Momentum & 0.9 \\
        Weight decay & 0.0 \\
        损失函数 & CrossEntropyLoss \\
        输入尺寸 & $128 \times 128$ \\
        输出类别 & 101 类，对应 0 到 100 岁 \\
        最优模型选择 & 验证集 MAE 最小时保存 \\
        \bottomrule
    \end{tabular}
\end{table}
    
  \item 在验证 MAE 下降时保存最佳模型权重到 \texttt{best\_age\_net.pth}；训练完成后使用最佳模型生成测试集预测写入 \texttt{custom.txt}。
\end{enumerate}

\section{实验内容}
本实验的主要任务为：基于给定网络结构与数据预处理，训练 AgeNet 实现年龄回归（通过 101 类分类后取期望作为连续年龄估计）。

\subsection{模型结构与关键实现}
关键模型定义摘自 \texttt{age\_net.py}：
\begin{lstlisting}
class AgeNet(nn.Module):
    def __init__(self, num_classes=101):
        super(AgeNet, self).__init__()
        self.features = nn.Sequential(
            nn.Conv2d(3,64,kernel_size=7,padding=3), nn.ReLU(), nn.MaxPool2d(2),
            nn.Conv2d(64,128,kernel_size=5,padding=2), nn.ReLU(),
            nn.Conv2d(128,128,kernel_size=5,padding=2), nn.ReLU(), nn.MaxPool2d(2),
            # 后续为多层 3x3 卷积与池化，最终通道 512
        )
        self.avgpool = nn.AdaptiveAvgPool2d((4,4))
        self.classifier = nn.Sequential(
            nn.Flatten(), nn.Linear(512*4*4,4096), nn.ReLU(), nn.Dropout(0.5),
            nn.Linear(4096,4096), nn.ReLU(), nn.Dropout(0.5),
            nn.Linear(4096,num_classes)
        )

    def forward(self,x):
        x = self.features(x)
        x = self.avgpool(x)
        x = self.classifier(x)
        return x
\end{lstlisting}

实现要点说明：
\begin{itemize}
  \item 使用 5 个卷积块逐步提取特征，保证训练时通道数与空间尺度与指导文档一致；
  \item 通过自适应池化确保进入全连接层前空间维为 $4\times4$，从而使线性层维度固定；
  \item 输出为 101 维 logits，训练时使用 \texttt{CrossEntropyLoss}，评价时通过 softmax 得到概率分布，再计算期望作为年龄预测：$\hat{age}=\sum_{i=0}^{100} p_i \cdot i$。
\end{itemize}

\subsection{训练流程（关键片段）}
训练脚本中的核心循环（节选自 \texttt{train\_ageNet.py}）：
\begin{lstlisting}
for epoch in range(EPOCH):
    model.train()
    running_loss = 0.0
    for i, (y,x) in enumerate(train_loader):
        x = x.to(device).float(); y = y.to(device).long()
        outputs = model(x)
        loss = criterion(outputs, y)
        optimizer.zero_grad(); loss.backward(); optimizer.step()
        running_loss += loss.item()

    # validation: 用概率分布求期望并计算 MAE
    with torch.no_grad():
        for i, (y,x) in enumerate(val_loader):
            outputs = model(x.to(device).float())
            probs = torch.softmax(outputs, dim=1)
            pred_age = (probs * ages).sum(dim=1)
            predsVal.append(pred_age.cpu().numpy())
            gtsVal.append(y.cpu().numpy())
    mae = np.mean(np.abs(predsVal - gtsVal))
    print(f'Epoch {epoch+1}/{EPOCH}  loss={running_loss:.4f}  val_mae={mae:.4f}')
    if mae < best_mae:
        best_mae = mae
        torch.save(model.state_dict(), best_path)
\end{lstlisting}

关键超参数：\texttt{EPOCH=100}, \texttt{TRAIN\_LR=0.001}, 优化器为 SGD（动量 0.9），无权重衰减。

\section{实验结果}
\subsection{训练日志摘录}
以下为训练过程中的关键输出片段（逐步记录验证 MAE 的变化与保存最佳模型信息）：
\begin{lstlisting}
Using device: cuda
Epoch 1/100  loss=1667.0581  val_mae=19.0350
Saved best model to /home/rin/Local/Github/XJTU-Labs/Computer-Vision/z-final-project/best_age_net.pth
...
Epoch 51/100  loss=1270.1084  val_mae=9.9518
Saved best model to /home/rin/Local/Github/XJTU-Labs/Computer-Vision/z-final-project/best_age_net.pth
...
Epoch 74/100  loss=1214.7470  val_mae=7.9926
Saved best model to /home/rin/Local/Github/XJTU-Labs/Computer-Vision/z-final-project/best_age_net.pth
...
Epoch 89/100  loss=1194.1046  val_mae=7.3164
Saved best model to /home/rin/Local/Github/XJTU-Labs/Computer-Vision/z-final-project/best_age_net.pth
...
Epoch 100/100  loss=1172.5933  val_mae=7.4276
=> training finished
Test results saved to custom.txt
\end{lstlisting}

说明：
\begin{itemize}
  \item "Using device: cuda" 表明训练时使用了 GPU，能显著加速卷积网络的训练；
  \item 训练早期（例如第 1 \textasciitilde{} 10 个 epoch）验证 MAE 快速从约 19 降到 11 \textasciitilde{} 12，表明网络从随机初始化开始快速学习到一定的年龄判别能力；
  \item 随后 MAE 持续缓慢下降，中后期（epoch 50 \textasciitilde{} 90）出现多个保存最佳模型的时刻，说明模型在验证集上不断找到更优的权重；
  \item 最佳记录出现在 epoch 89（val\_mae = 7.3164），训练结束时最终验证 MAE 为 7.4276，整体表现稳定在约 7.3 \textasciitilde{} 7.9 范围；
  \item 每次出现 "Saved best model to ..." 对应验证 MAE 刷新，这里保存路径为项目根下的 \texttt{best\_age\_net.pth}。
\end{itemize}

\subsection{结果分析}
\begin{itemize}
  \item 从数值上看，最终验证 MAE 约为 7.3 \textasciitilde{} 7.4，这是从零初始化、基于 VGG 风格小型网络在 APPA-REAL 上能取得的合理基线精度；
  \item 模型以分类 + 期望回归的方式输出连续年龄估计：相比直接回归，分类后期望的做法可以通过 softmax 建模不确定性并稳定训练；
  \item 训练日志中 loss 为每个 epoch 的累加和（注意不是平均损失），因此与常见的每 batch 平均 loss 值在数量级上不同；
  \item 可进一步改进点：使用更强的数据增强、学习率调度（如 StepLR 或 CosineAnnealing）、预训练骨干（如 ImageNet 预训练的 VGG/ResNet）以及更大的批量或长时间训练。 
\end{itemize}

\section{总结}
本实验按照指导实现并训练了 \texttt{AgeNet} 模型，验证集上最终 MAE 约为 7.3，训练过程中多次记录并保存了验证集上的最优权重，测试预测结果写入 \texttt{custom.txt}。实现遵循了实验要求：固定输入尺寸、5 个卷积块、3 个全连接层、SGD 优化及 100 个 epoch 上限。

\section{附录：主要源码与运行命令}

\subsection{运行命令}
\begin{verbatim}
# 在准备好 DATASET/ 后运行：
python train_ageNet.py
\end{verbatim}

\subsection{主要源码}

\subsubsection{\texttt{age\_net.py}}

\begin{lstlisting}
import torch
import torch.nn as nn


class AgeNet(nn.Module):
    """AgeNet classification model producing 101-way logits (ages 0..100).

    Designed for 128x128 RGB input following the provided architecture.
    """
    def __init__(self, num_classes=101):
        super(AgeNet, self).__init__()

        self.features = nn.Sequential(
            nn.Conv2d(3, 64, kernel_size=7, stride=1, padding=3),
            nn.ReLU(inplace=True),
            nn.MaxPool2d(kernel_size=2, stride=2),

            nn.Conv2d(64, 128, kernel_size=5, padding=2),
            nn.ReLU(inplace=True),
            nn.Conv2d(128, 128, kernel_size=5, padding=2),
            nn.ReLU(inplace=True),
            nn.MaxPool2d(kernel_size=2, stride=2),

            nn.Conv2d(128, 256, kernel_size=3, padding=1),
            nn.ReLU(inplace=True),
            nn.Conv2d(256, 256, kernel_size=3, padding=1),
            nn.ReLU(inplace=True),
            nn.Conv2d(256, 256, kernel_size=3, padding=1),
            nn.ReLU(inplace=True),
            nn.MaxPool2d(kernel_size=2, stride=2),

            nn.Conv2d(256, 512, kernel_size=3, padding=1),
            nn.ReLU(inplace=True),
            nn.Conv2d(512, 512, kernel_size=3, padding=1),
            nn.ReLU(inplace=True),
            nn.Conv2d(512, 512, kernel_size=3, padding=1),
            nn.ReLU(inplace=True),
            nn.MaxPool2d(kernel_size=2, stride=2),

            nn.Conv2d(512, 512, kernel_size=3, padding=1),
            nn.ReLU(inplace=True),
            nn.Conv2d(512, 512, kernel_size=3, padding=1),
            nn.ReLU(inplace=True),
            nn.Conv2d(512, 512, kernel_size=3, padding=1),
            nn.ReLU(inplace=True),
            nn.MaxPool2d(kernel_size=2, stride=2),
        )

        # ensure spatial size is 4x4 before the first FC
        self.avgpool = nn.AdaptiveAvgPool2d((4, 4))

        self.classifier = nn.Sequential(
            nn.Flatten(),
            nn.Linear(512 * 4 * 4, 4096),
            nn.ReLU(inplace=True),
            nn.Dropout(0.5),
            nn.Linear(4096, 4096),
            nn.ReLU(inplace=True),
            nn.Dropout(0.5),
            nn.Linear(4096, num_classes),
        )

        self._initialize_weights()

    def forward(self, x):
        x = self.features(x)
        x = self.avgpool(x)
        x = self.classifier(x)
        return x

    def _initialize_weights(self):
        for m in self.modules():
            if isinstance(m, nn.Conv2d):
                nn.init.kaiming_normal_(m.weight, mode='fan_out', nonlinearity='relu')
                if m.bias is not None:
                    nn.init.constant_(m.bias, 0)
            elif isinstance(m, nn.BatchNorm2d):
                nn.init.constant_(m.weight, 1)
                nn.init.constant_(m.bias, 0)
            elif isinstance(m, nn.Linear):
                nn.init.normal_(m.weight, 0, 0.01)
                nn.init.constant_(m.bias, 0)
\end{lstlisting}

\subsubsection{\texttt{train\_ageNet.py}}

\begin{lstlisting}
import os
import numpy as np
import torch
import torch.nn as nn
from torch.utils.data import DataLoader

from helperT import get_img_dataloaders, test_cel
from age_net import AgeNet


def train_ageNet(base_dir='DATASET/'):
    # hyper-parameters
    EPOCH = 100
    TRAIN_LR = 0.001
    MOMENTUM = 0.9
    WEIGHT_DECAY = 0.0

    # data
    train_loader, val_loader, test_loader = get_img_dataloaders(base_dir)

    device = torch.device('cuda' if torch.cuda.is_available() else 'cpu')
    print('Using device:', device)

    model = AgeNet(num_classes=101).to(device)
    optimizer = torch.optim.SGD(model.parameters(), lr=TRAIN_LR, momentum=MOMENTUM, weight_decay=WEIGHT_DECAY)
    criterion = nn.CrossEntropyLoss()

    best_mae = 1e9
    best_path = os.path.join(os.path.dirname(__file__), 'best_age_net.pth')

    ages = torch.arange(0, 101).float().to(device)

    for epoch in range(EPOCH):
        model.train()
        running_loss = 0.0
        for i, (y, x) in enumerate(train_loader):
            x = x.to(device).float()
            y = y.to(device).long()

            outputs = model(x)
            loss = criterion(outputs, y)

            optimizer.zero_grad()
            loss.backward()
            optimizer.step()

            running_loss += loss.item()

        # validation
        model.eval()
        predsVal = []
        gtsVal = []
        with torch.no_grad():
            for i, (y, x) in enumerate(val_loader):
                x = x.to(device).float()
                y = y.to(device).long()
                outputs = model(x)
                probs = torch.softmax(outputs, dim=1)
                pred_age = (probs * ages).sum(dim=1)
                predsVal.append(pred_age.cpu().numpy())
                gtsVal.append(y.cpu().numpy())

        predsVal = np.concatenate(predsVal, axis=0)
        gtsVal = np.concatenate(gtsVal, axis=0)
        mae = np.mean(np.abs(predsVal - gtsVal))

        print(f'Epoch {epoch+1}/{EPOCH}  loss={running_loss:.4f}  val_mae={mae:.4f}')

        if mae < best_mae:
            best_mae = mae
            torch.save(model.state_dict(), best_path)
            print('Saved best model to', best_path)

    print('=> training finished')
    return best_path, test_loader


if __name__ == '__main__':
    best_path, test_loader = train_ageNet('DATASET/')

    # Load best model and run test to produce predictions file
    device = torch.device('cuda' if torch.cuda.is_available() else 'cpu')
    model = AgeNet(num_classes=101).to(device)
    model.load_state_dict(torch.load(best_path, map_location=device))
    model.to(device)

    # helperT.test_cel expects the model to be on cuda and will call .cuda() on inputs.
    # If running on CPU, test_cel will still produce outputs if model and inputs are on CPU.
    prediction = test_cel(model, test_loader, os.path.join(os.path.dirname(__file__), 'custom.txt'))
    print('Test results saved to custom.txt')

\end{lstlisting}

\end{document}
